\chapter{Các nội dung tính toán}
    \section{Thông lượng dữ liệu và hiệu suất băng thông}
 
        Thông lượng ứng dụng thực tế $R_{app}$ là lượng dữ liệu hữu ích (payload) còn lại sau khi trừ overhead giao thức tại tất cả các lớp. Hiệu suất băng thông được định nghĩa:
        \[ \eta = \frac{R_{app}}{R_{PHY}} \]
        
        \textbf{Zigbee (E18-ZG120 / CC2530).} Tốc độ vật lý IEEE 802.15.4 tại băng 2,4~GHz đạt $R_{PHY} = 250\,\text{kbps}$ \cite{ieee802154_2020}. Mỗi frame Zigbee mang tối đa 99 byte payload ứng dụng sau khi trừ 28 byte overhead của các lớp MAC/NWK/APS trên tổng PSDU 127 byte \cite{zigbee_spec_r21}. Một chu kỳ phát--nhận ACK bao gồm: thời gian truyền frame ($127\,\text{B} \times 8 / 250\,\text{kbps} = 4{.}064\,\text{ms}$), khoảng LIFS ($40\,\text{sym} \times 16\,\mu\text{s} = 0{.}640\,\text{ms}$) và frame ACK ($0{.}544\,\text{ms}$), tổng $T_{air} \approx 5{.}25\,\text{ms}$.
        
        Nút thắt thực tế không phải là air interface mà là giao tiếp UART từ module E18-ZG120 lên LAN-MCU. Module đóng gói mỗi khung Zigbee (127 byte PSDU) vào khung HEX của Ebyte với 19 byte framing overhead, tạo ra frame UART 146 byte. Tại tốc độ $115200\,\text{bps}$ với 10 bit/byte (8N1), thời gian truyền một frame là:
        \[ T_{UART,Zigbee} = \frac{146 \times 10}{115200} \approx 12{.}67\,\text{ms} \]
        Điều này khiến UART chậm hơn air interface $12{.}67\,\text{ms} / 5{.}25\,\text{ms} \approx 2{.}4\times$, và trở thành điểm nghẽn của toàn chặng Zigbee.
        
        Trong kiểm thử thực tế với lệnh \texttt{reportTemperature()} -- chỉ mang 2--4 byte dữ liệu hữu ích nhưng phải gánh toàn bộ overhead ZCL -- thông lượng ứng dụng chỉ đạt $R_{app} \approx 8{.}2\,\text{kbps}$. Khi nhiều node phát đồng thời, cơ chế CSMA/CA ép các node lùi thời gian phát ngẫu nhiên (random backoff) khiến tổng thông lượng không tăng mà còn bị bão hòa tại ngưỡng này; vượt qua đó hàng đợi APS sẽ tràn và xảy ra mất gói (packet loss).
        
        \textbf{LoRa (Wio-E5 / STM32WLE5JC) -- Chế độ TEST/P2P.} Hệ thống sử dụng chế độ \texttt{AT+MODE=TEST} thay vì LoRaWAN, cho phép bỏ qua hoàn toàn giới hạn duty cycle 1\% của ETSI và stack LoRaWAN; thông lượng chỉ bị giới hạn bởi Time-on-Air. Tốc độ PHY tại SF7/BW125\,kHz đạt $R_{PHY} = 5{.}47\,\text{kbps}$ \cite{semtech_lora_design_guide}. Với gói 50 byte, Time-on-Air tính theo công thức Semtech là $T_{air} \approx 97{.}5\,\text{ms}$. Mỗi chu kỳ phát còn bao gồm thời gian xử lý AT command và nhận phản hồi \texttt{+TEST: TX DONE} ($\approx 16{.}1\,\text{ms}$), tổng $T_{cycle} \approx 113{.}6\,\text{ms}$, dẫn đến:
        \[ R_{app,LoRa} = \frac{50 \times 8}{113{.}6\,\text{ms}} \approx 3{.}52\,\text{kbps} \]
        Giới hạn payload tối đa qua AT command là 169 byte/gói do độ dài lệnh tối đa 528 ký tự \cite{seeed_wioe5_at_cmd}.
        
        \textbf{BLE (Native ESP32-S3).} BLE được triển khai trực tiếp trên LAN-MCU nên không có tầng UART trung gian. Thông lượng bị giới hạn bởi Connection Interval ($CI = 20\,\text{ms}$) và kích thước ATT payload. Ở cấu hình mặc định ATT MTU = 23 byte (20 byte dữ liệu), thông lượng chỉ đạt $\approx 8\,\text{kbps}$ \cite{bluetooth_core_v5}. Khi kích hoạt Data Length Extension (DLE) và đàm phán MTU lên 247 byte (244 byte ATT payload), thông lượng cải thiện đáng kể:
        \[ R_{app,BLE,DLE} = \frac{244 \times 8}{20\,\text{ms}} = 97{.}6\,\text{kbps} \]
    